%
% This document is available under the Creative Commons Attribution-ShareAlike
% License; additional terms may apply. See
%   * http://creativecommons.org/licenses/by-sa/3.0/
%   * http://creativecommons.org/licenses/by-sa/3.0/legalcode
%
% Copyright 2010 Jérôme Pouiller <jezz@sysmic.org>
%

% Configuration de Beamer

\usetheme{Warsaw}

% Quiris
%\usecolortheme[RGB={238,127,0}]{structure}
% Apollo
%\usecolortheme[RGB={0,56,126}]{structure}
% Sysmic
\usecolortheme[RGB={181,0,0}]{structure}

% Headers et footers
\useoutertheme[subsection=false]{smoothbars}
% Rectangle dans les itemize
\useinnertheme{rectangles}
% Pas de symboles pour la navigation
\setbeamertemplate{navigation symbols}{}

% Logo en filigranne
\setbeamertemplate{background canvas}{
  \tikz {
    \node at (0,0) {};
    \node[inner sep=0pt, opacity=0.2] at (0.5\paperwidth,-0.8\paperheight) {\includegraphics[height=23.4mm,width=128mm]{pics/filigranne-sysmic}};
  }
}

% Ajout des numéros de pages
\newcommand*\oldinsertshortitle{}%
\let\oldinsertshorttitle\insertshorttitle%
\renewcommand*\insertshorttitle{\oldinsertshorttitle\hfill\insertframenumber\,/\,\inserttotalframenumber}

\usepackage{ucs}              % La sortie est en UTF8
\usepackage[utf8x]{inputenc}  % L'entrée aussi
\usepackage[T1]{fontenc}      % Pour la césure des mots accentués
\usepackage{lmodern}          % Pour les lettres francaises
\usepackage[french]{babel}    % Pour la sortie francaise
\usepackage{times}
\usepackage{mathptmx}         % Font (don't forget to install latex-font-recommended)

% This document is available under the Creative Commons Attribution-ShareAlike
% License; additional terms may apply. See
%   * http://creativecommons.org/licenses/by-sa/3.0/
%   * http://creativecommons.org/licenses/by-sa/3.0/legalcode
%
% Copyright 2010 Jérôme Pouiller <jezz@sysmic.org>
%

\usepackage{tikz}
\usepackage{pgfplots}
\usepackage{tikz-timing}
\usetikzlibrary{patterns}
\usetikzlibrary{shapes}
\usetikzlibrary{snakes}
\usetikzlibrary{positioning}
\usetikzlibrary{calc}

\tikzstyle{cgreen}  = [fill=green!20!white,  draw=green!50!black,  line width=1pt, rounded corners=1pt]
\tikzstyle{cred}    = [fill=red!20!white,    draw=red!50!black,    line width=1pt, rounded corners=1pt]
\tikzstyle{cblue}   = [fill=blue!20!white,   draw=blue!50!black,   line width=1pt, rounded corners=1pt]
\tikzstyle{cpurple} = [fill=purple!20!white, draw=purple!50!black, line width=1pt, rounded corners=1pt]
\tikzstyle{corange} = [fill=orange!20!white, draw=orange!50!black, line width=1pt, rounded corners=1pt]
\tikzstyle{ccyan}   = [fill=cyan!20!white,   draw=cyan!50!black,   line width=1pt, rounded corners=1pt]
\tikzstyle{cbrown}  = [fill=brown!20!white,  draw=brown!50!black,  line width=1pt, rounded corners=1pt]
\tikzstyle{cyellow} = [fill=yellow!20!white, draw=yellow!50!black, line width=1pt, rounded corners=1pt]

\def\timeline#1#2#3{
  \draw[xstep=1,ystep=1.5,gray,very thin] (0,0) grid (#1.5,#2);
  \draw[->, line width=1pt] (0,#2) -- (#1.5,#2) coordinate (x axis);
  \draw[line width=1pt] (0,#2) -- (0,0) coordinate (x axis);
  \foreach \x in {0,1,...,#1}
    \draw (\x.5,#2) node[anchor=north] {\tiny\x};
  \foreach \y/\ytext in {#3}
    \draw (0,\y) node[anchor=east] {\ytext};
}

\def\lo#1{ -- ++(#1,0) }
\def\hi#1{ |- ++(#1,1) -- ++(0,-1) }
\def\hop#1#2{\lo{#1}\hi{#2}}

% Period End
\def\pe#1#2#3{
   \begin{scope}
     \clip  (#1 - 0.2, #2 - 0.2) rectangle (#1, #2 + 0.2);
     \draw[#3] (#1, #2) circle (3pt);
   \end{scope}
}

% Period Begin
\def\pb#1#2#3{
   \begin{scope}
     \clip  (#1 + 0.2, #2 - 0.2) rectangle (#1, #2 + 0.2);
     \draw[#3] (#1, #2) circle (3pt);
   \end{scope}
}

% Period Begin and End
\def\p#1#2#3{
   \draw[#3] (#1, #2) circle (3pt);
}



\ifpdf
  \usepackage{embedfile} 
\else
  \newcommand{\embedfile}[1]{}
\fi

% Produce errors like:
% main.tex:54: Undefined control sequence.
% \@EveryShipout@Output ...EveryShipout@Org@Shipout
%                                                   \box \@cclv
%\usepackage[sections,displaymath]{preview}
%\PreviewEnvironment{tikzpicture}
%\PreviewEnvironment{center}
%\PreviewEnvironment{frame}

\usepackage{hyperref} % Doit être chargé avant ntheorem
\newcommand{\email}[1]{\href{mailto:#1}{\nolinkurl{<#1>}}}
\newcommand{\man}[1]{\emph{#1}}
\newcommand{\file}[1]{\lstinline[backgroundcolor=\color{{rgb}{1,1,0.8}},language=]{#1}}
\newcommand{\cmd}[1]{\lstinline[backgroundcolor=\color{{rgb}{1,1,0.8}},language=]{#1}}
\renewcommand{\c}[1]{\lstinline[backgroundcolor=\color{{rgb}{1,1,0.8}},language=c]{#1}}

\usepackage[newcommand]{ragged2e}

% \usepackage[]{version} 
% \usepackage{ntheorem} 
% \setlength{\theorempreskipamount}{0.8ex plus 0.9ex minus 0.1ex}
% \setlength{\theorempostskipamount}{0.8ex plus 0.9ex minus 0.1ex}
% \theoremprework{\vspace{3mm}\hrule}
% \theorempostwork{\hrule\vspace{3mm}}
% \theorembodyfont{\itshape}
% \theoremseparator{.}
% \newtheorem{quest}{Question}[section]

% \theorembodyfont{\normalfont}
% \theoremseparator{}
% \theoremstyle{break}
% \newtheorem*{ans}{Réponse}

% \setlength{\theoremindent}{3mm}
% \theoremseparator{:}
% \theorembodyfont{\itshape}
% \theoremstyle{plain}
% \newtheorem*{man}{Documentation utile}
% \theorembodyfont{\normalfont}
% \newtheorem*{hint}{Remarque}
% \newtheorem*{note}{Note pédagogique}

% \usepackage[vmargin=25mm,hmargin=15mm]{geometry}          % Marges peronnalisées

% Quelques règlage de mise en page
%\renewcommand{\baselinestretch}{1.2} % taille de l'interligne
\setlength{\parindent}{0pt}
\setlength{\parskip}{0.9ex plus 0.5ex minus 0.2ex}

%\usepackage{fancyhdr}        % Fancy page headers
%\lhead{}                     % Top-Left
%\chead{}                     % Top-Center
%\rhead{}                     % Top-Right
%\lfoot{}                     % Bottom-Left
%\cfoot{}                     % Bottom-Center
%\rfoot{}                     % Bottom-Right
%\pagestyle{fancy}

\usepackage{listings}         % Pour mettre en page du code source
\usepackage{color}            % Pour les lien en couleur dans le pdf
\definecolor{colBg}        {rgb}{1,1,0.8}
\definecolor{colKeys}      {rgb}{0,0,1}
\definecolor{colComments}  {rgb}{1,0,0}
\definecolor{colString}    {rgb}{0,0.5,0}
\definecolor{colBasic}     {rgb}{0,0,0}
\definecolor{colIdentifier}{rgb}{0,0,0}
\lstset{%                     % Basic style
  numbers=left,%
  stepnumber=10,%
  numberstyle=\scriptsize,%
%
  basicstyle=\ttfamily\normalsize\color{colBasic},%
  commentstyle=\normalsize\itshape\color{colComments},%
  identifierstyle=\color{colIdentifier},%
  keywordstyle=\bf\ttfamily\color{colKeys},%
  stringstyle=\color{colString},%
  backgroundcolor=\color{colBg},%
%
%  mathescape=true,%
  extendedchars=false,%
%  tabsize=4,%
  columns=flexible,%
  fontadjust=true,%
  frame=lines,%
  showspaces=false,%
  showstringspaces=false,%
%
  emptylines=1,%
  breaklines=true,%
  breakautoindent=true,%     % Inutile avec breaklines=false
  literate={é}{{\'e}}1 {è}{{\`e}}1 {ô}{{\^o}}1 {à}{{\`a}}1 {ç}{{\c{c}}}1 
}
\lstset{language=C++}        % ... En C++

\definecolor{darkgreen}{rgb}{0, 0.7, 0}
\definecolor{darkgreen2}{rgb}{0, 0.5, 0}
\definecolor{red2}{rgb}{0.8, 0, 0}
\lstdefinelanguage{diff} {
    morecomment=[f][\color{darkgreen}][0]{+},
    morecomment=[f][\color{red}][0]{-},
    morecomment=[f][\itshape\color{darkgreen2}][0]{+++},
    morecomment=[f][\itshape\color{red2}][0]{---},
    moredelim=[l][\color{cyan}]{\ \@\@},
    moredelim=*[l][\color{blue}]{\@\@},
}

\hypersetup{colorlinks=true,plainpages=false,urlcolor=blue,linkcolor=}






